\section{Discussion}

\subsection{Geometry and noncommuting limits}

Within the controlled weak-cubic family and the declared local annulus, the
essential change is geometric rather than a uniform frequency displacement.
The local Morse--Bott ring unfolds into isolated Morse minima and saddles.
The isotropic ZGV critical ring is Morse--Bott: radial localization is
nondegenerate, but continuous rotation leaves a tangential Hessian zero mode
(Theorem~\ref{thm:morse-bott-ring}).  Cubic anisotropy lifts that zero mode and
organizes the local stationary set into alternating Cartesian Morse points
(Theorem~\ref{thm:cubic-morse-splitting}).  The extra nondegenerate direction
adds angular stationary phase to the existing radial stationary phase.  This
explains why the ring and point regimes carry different time exponents without
describing the change as a phase transition.

The two exponents belong to noncommuting limits.  Write
\(\alpha=\varepsilon V_4\).  The ring-scale response is proportional to
\(t^{-1/2}J_0(\alpha t)\), whereas the isolated-point scale in the controlled
overlap is \(\lvert\alpha\rvert^{-1/2}t^{-1}\).  The latter prefactor follows
from the tangential Hessian curvature in
Eq.~\eqref{eq:overlap-hessian-signatures}; point multiplicity cancels the
fourfold curvature factor but not \(\lvert\alpha\rvert^{-1/2}\).  Its growth as
\(\alpha\to0\) signals loss of uniformity when the points coalesce back into a
ring, not divergence of the physical solution.  At
\(t\sim\lvert\alpha\rvert^{-1}\), both asymptotic scales are of order
\(\lvert\alpha\rvert^{1/2}\), as summarized in
Eq.~\eqref{eq:noncommuting-green-limits}.

The joint-limit Bessel law and fixed-anisotropy exact-Morse sum are therefore
complementary.  The
uniform Bessel theorem controls compact \(\tau=\alpha t\)
(Theorem~\ref{thm:uniform-bessel-crossover}); its large-parameter edge is
matched only through the separately proved growing-\(\lvert\tau\rvert\)
overlap and Eq.~\eqref{eq:bessel-overlap-response}.  At fixed nonzero
anisotropy, the exact-Morse theorem instead uses full critical frequencies,
amplitudes, and Cartesian Hessians (Theorem~\ref{thm:fixed-anisotropy-decay}).
Its remainder is not uniform as \(\varepsilon\to0\), while a first-order
Bessel phase is not indefinitely accurate at fixed \(\varepsilon\).  Neither
formula can be substituted for the other outside the regime it controls.

\subsection{Relation to prior work}

Within the scope of this comparison, both endpoint geometries are established
results rather than discoveries of the present study.  The isotropic ZGV
critical ring and its \(t^{-1/2}\) response follow the classical Lamb-wave
setting \citep{lamb1917waves,prada2008local,prada2008power}.  Directional
anisotropic shifts were also established \citep{prada2009anisotropy}, and
isolated anisotropic ZGV points were previously reported.  In particular, the
four minima and four saddles were previously reported together with their
interference, and beating responses were already reported for cubic silicon
\citep{kiefer2023beating,kiefer2025skewing}.  Full-spectrum algorithms for
locating such points are likewise prior infrastructure
\citep{kiefer2023computing}.

The distinction of the present calculation is the controlled,
coefficient-resolved link between those known endpoints.  Morse--Bott
classification and perturbation are established mathematics
\citep{bott1954critical,banyaga2013cascades}.  Here transverse nondegeneracy is
verified on the exact Lamb determinant and independently on the tracked
elastic branch.  The fourfold coefficient is obtained from the full
generalized-eigenvalue sensitivity, and the radial coefficient retains the
differentiated-mode contribution.  Finite differences, compensated radial
shifts, remainder scaling, sign reversal, Cartesian Hessian inertia, and
annular index closure then test the reduction.  Thus the local four-plus-four
result follows from computed coefficients and controlled remainders, not from
point-group symmetry alone or from fitting the final critical locations.

The asymptotic tools also have clear precedents.  Broken-symmetry uniform
approximations with Bessel modulation connect continuous families and isolated
objects in other wave settings \citep{creagh1996broken,brack1999uniform,
brack2009closed}.  Two-dimensional stationary phase is standard and has been
used for anisotropic guided-wave radiation
\citep{velichko2007excitation,chapuis2010focusing,karmazin2013caustics}.
Accordingly, the contribution is neither a new Bessel identity nor a new
stationary-phase theorem.  It is their fully normalized Lamb-wave realization
with an elasticity-derived Bessel argument, an explicit phase-validity window,
and a constant-and-phase match to the first-order Morse overlap.  Separately,
the fixed-anisotropy endpoint is compared numerically with the exact-Hessian
point sum without fitted amplitude, phase, frequency, or time shifts.

\subsection{Scope and boundary statements}

The canonical Bessel law holds only within the specified nonnodal source and
angular-weight class.  A general source and its angular weight jointly produce
the Fourier--Bessel series of Eq.~\eqref{eq:weighted-bessel-series}; the leading
kernel need not reduce to a single \(J_0\) term.  If the ring amplitude
vanishes, or if a detector is nodal at a particular Morse point, the leading
contribution changes as described in Supplementary
Section~\ref{rem:nodal-channels}.
Cancellation times do not alter the generic exponent, but they do make a
pointwise relative error or local slope ill conditioned.  The declared
source-radius and window sweeps support robustness only over their declared
ranges, not source independence.

The generality of the result has two levels.  At theorem level, any simple,
uniformly gapped finite-radius branch with nonzero radial curvature and a
dominant nonzero fourfold perturbation coefficient has the same local
Morse--Bott-to-Morse structure and exact sufficiently-small four-plus-four
count under the stated remainder hypotheses.  A pure leading fourfold phase,
together with the declared nonnodal source and phase-remainder conditions, is
additionally required for the single-\(J_0\) law and its matched asymptotic
regimes; generic higher harmonics change the angular canonical integral.  At
instance level, coefficient values,
finite-resolution roots, and response comparisons are demonstrated for one
selected Lamb branch and one cubic perturbation family.  The silicon stress
test assesses only the search and classification workflow.

The geometric count is also restricted to one simple branch in the declared
local annulus.  Branch contamination is controlled here by tracking and the
eigengap gate of Eq.~\eqref{eq:eigengap-rejection}; if the eigengap closes, a
scalar sensitivity and reduced resolvent no longer suffice.  Degenerate,
matrix-valued perturbation theory would then be required.  Likewise, zero
radial curvature, a zero-wavenumber stationary point, or a singular exact
Hessian would demand a different canonical scaling.  The computed local count
does not enumerate every Lamb or shear-horizontal branch.

Fourfold symmetry alone does not fix the number of points.  The pure leading
fourfold term is a result of the declared first-order cubic family.  If the
coefficient \(V_4\) is zero, higher harmonics or the second-order effective
angular potential can control the splitting; competing higher harmonics can
also change the root count.  The coefficient \(B\) is not part of the
controlled remainder: the explicit \(\varepsilon qB\) term controls the radial
shift, and \(B\) enters the phase-error and overlap budget through the effective
second-order phase.  Omitting \(B\) therefore changes both the critical
locations and temporal validity.  These cases require a new
coefficient audit rather than extrapolation of the present four-plus-four
formula.

Neither temporal approximation is indefinitely phase accurate outside its
declared late-time window.  Compact-\(\tau\) collapse, the growing-parameter
overlap, and the fixed-anisotropy comparison impose different error
conditions.  The accumulated inter-resolution phase-discrepancy estimator in
Eq.~\eqref{eq:phase-error-certificate} restricts the accepted numerical time
window but does not bound the unknown continuum phase error.  The operational
crossover time in Fig.~\ref{fig:decay-crossover} is found inside a
theory-centred bracket and is therefore a consistency diagnostic, not
independent evidence for inverse-rate scaling; the complex Bessel collapse is
the primary numerical test.  Critical-frequency features in
the lossless continuum should not be interpreted as discrete resonances, and
their exact separation remains twice the defined modulation magnitude.

Finally, the present model is a linear, infinite, homogeneous, lossless plate.
Damping generally introduces complex dispersion and additional decay; finite
boundaries introduce reflections and a discrete modal spectrum; layers,
interfaces, and defects add branches, scattering, and new angular weights.
Each extension requires renewed branch
projection, coefficients, and phase-error analysis.  The finite-anisotropy
silicon stress test is only a robustness check on search, classification, and
boundary consistency checking.  It supplies no weak-parameter theorem and no
global point count.  This computation-only study includes no experiments.

\subsection{Nonlinear amplitude equations are a separate problem}

The normal form used here is only a linear spectral reduction, not a nonlinear
amplitude equation.  Its frequency sensitivities and dispersive exponents
describe propagation after an impulsive load under the linear constitutive
model.  They do not contain nonlinear coupling, saturation, drive-dependent
detuning, mode competition, or pattern selection.  In particular, the local
Morse geometry can organize the linear spectral input to a later theory, but
it does not determine nonlinear states or their stability.

Nonlinear amplitude equations are therefore a separate problem.  Such a study
would require a nonlinear constitutive law, nonlinear traction conditions, a
multiscale solvability calculation, and declared scalings for forcing,
damping, and detuning.  Resonant projections and coupling coefficients would
have to be derived independently, followed by stability analysis.  None of
those equations or coefficients is obtained from the present Green response,
and no nonlinear saturation or pattern-selection result is inferred here.
